\subsection{Owned choice, dissemilattices, and the core}
\label{app:core-basis}

The core contains a classical variety that gives a finite basis for
another fragment and excludes a finite residual bound for the core itself.
Write $+$ for union and $xy$ for $x\star_a y$ in this subsection.
A \emph{meet-distributive bisemilattice}, also called a
\emph{dissemilattice}, satisfies the three semilattice laws for each
operation and the one distributive law
\begin{equation}\label{core:eq:dissemilattice}
 x(y+z)=xy+xz.
\end{equation}
This is a basis of seven identities for that class.

\begin{proposition}[A seven-equation basis]\label{core:prop:choice-union-basis}
For every $K\geq3$ and owner $a$, the winner-set algebra in the
signature $\{\star_a,\cup\}$ generates exactly the variety of
meet-distributive bisemilattices. Its identities are therefore
axiomatized by the seven equations just described.
\end{proposition}
\begin{proof}
Both operations are semilattices. For
\eqref{core:eq:dissemilattice}, if $x=\{a\}$, both sides are
$\{a\}$. Otherwise, if both $y,z$ equal $\{a\}$, both sides
again equal $\{a\}$. If neither equals $\{a\}$, both sides are
$x\cup y\cup z$. If exactly one equals $\{a\}$, the same union
is obtained: the absorbed branch on the right contributes the
owner singleton and the other branch contributes $x$.

Choose distinct $a,b,c$ and take the three-element subalgebra
\[
 B=\{a\},\qquad N=\{a,b\},\qquad P=\{a,b,c\}.
\]
Union is maximum in $B<N<P$; owned choice is maximum in
$N<P<B$. This is the classical generator $3_m$ of the
meet-distributive bisemilattices
\cite[Section~3.3, Theorem~3.9 and the following paragraph]{HardingRomanowska2017}.
In particular it fails the other distributive law:
\[
 N+(BP)=N,\qquad (N+B)(N+P)=P.
\]
The winner-set algebra belongs to that variety and contains its
generator, so it generates precisely that variety.
\end{proof}

This is the raw owned-choice fragment, distinct from the guarded
one-owner theory of \cref{owner:thm:finite-basis}. Its equation
problem is coNP-complete by \cref{thm:raw-pairs}. A finite basis
does not impose a polynomial bound on equational derivations.

\begin{proposition}[Dissemilattices inside the core]
\label{core:prop:dissemilattice-lift}
Fix an owner $a$. Every meet-distributive bisemilattice $(D,+,\cdot)$
expands to a member of $\operatorname V(\mathbf C)$ by setting
\[
 x\star_a y=x\cdot y,\qquad
 x\star_r y=x+y\quad(r\ne a).
\]
This expansion preserves the congruence lattice. The variety
$\operatorname V(\mathbf C)$ has no finite residual bound.
\end{proposition}
\begin{proof}
The same three winner sets $B,N,P$ form a subalgebra for every
owned choice. For $r\ne a$ none is the singleton $\{r\}$, so
$\star_r$ is union on them. Thus this subalgebra is exactly the
displayed expansion of $3_m$.

Duplicating operation symbols commutes with products, subalgebras,
and homomorphic images. Since $3_m$ generates all meet-distributive
bisemilattices, their displayed expansions belong to
$\operatorname V(\mathbf C)$. Both original operations occur in
the expansion, and every added operation repeats one of them,
so its congruences are unchanged.

There are infinitely many nonisomorphic subdirectly irreducible
meet-distributive bisemilattices
\cite[Section~3.3, following Theorem~3.9]{HardingRomanowska2017}.
Their expansions remain subdirectly irreducible and pairwise
nonisomorphic. A finite signature has only finitely many algebras
of any bounded finite cardinality, up to isomorphism. Consequently
no finite cardinal bounds the subdirectly irreducible core models.
\end{proof}

This excludes the finite-residual-bound hypothesis of Willard's
theorem for the core. It does not prove nonfinite basability: the
embedded dissemilattice variety itself has the seven-equation basis.
Nor does it settle the residual bound after adding winner constants
or internal role switching; the displayed expansions supply neither.

The generic associativity and distributivity calculations, and the
three-state embedding for three roles, are kernel-checked in
\path{lean/CoreModels.lean}. The dissemilattice generation and
irreducible-family results are imported from the cited literature.
The expansion and congruence arguments above are written proofs.
\path{checks/check_core_basis.py} checks the seven sound identities,
the concrete embeddings, and the failed reverse distributive law
independently with winner sets.
